% GENERATED FILE. Edit canonical lessons, then run npm run generate:books.
% canonical-source-hash: fnv1a64:b0131dedf84c78fe
% canonical-lessons: MW-C189-thabho, MW-C189-it, MW-C189-chuno, MW-C189-dehli, MW-C189-chaubaro

\chapter{Walls and Doorways}
\label{ch:mw-walls-and-doorways}
\hlchaptermodality{189}

\begin{chapteropening}
\textbf{By the end of this chapter:} \emph{I can name what a house is built from.}

Complete the last lesson of chapter 189: I can name what a house is built from.
\end{chapteropening}

\section[\texorpdfstring{\mw{थांभो} (thā̃bho)}{थांभो (thā̃bho)}]{\mw{थांभो} (thā̃bho) — a pillar}

\label{lesson:MW-C189-thabho}



\begin{quote}
Before the new one: say the Marwadi for a hut, then the Marwadi for a thatched shelter.
\end{quote}

\subsection*{You'll want to know: \mw{थांभो}}
\textbf{\mw{थांभो}} — \emph{thā̃bho} — \textquotedblleft{}a pillar\textquotedblright{}.

\begin{quote}
\mw{थांभो} \mw{लांबो} \mw{है।}
\end{quote}

\begin{quote}
\emph{thā̃bho lā̃bo hai.}
\end{quote}

\begin{quote}
The pillar is tall.
\end{quote}

A masculine noun.

\begin{cousinweb}
\textbf{\mw{थांभो}} (\emph{thā̃bho}) is inherited from Sanskrit \emph{stambha}, \textquotedblleft{}a pillar, a post\textquotedblright{}; CDIAL lists Old Marwari \emph{thāma}, \textquotedblleft{}a pillar\textquotedblright{} (CDIAL 13682).
\end{cousinweb}

\begin{grammarlens}[title={What you've built}]
One more word for this chapter.
\end{grammarlens}

\subsection*{Your turn}
\begin{itemize}
\raggedright
  \item \emph{Say it:} \emph{thā̃bho}
  \item \emph{Say it:} \emph{thā̃bho}, once more
  \item \emph{Say it:} \emph{thā̃bho lā̃bo hai.} — The pillar is tall.
  \item \emph{Recall:} say the Marwadi for a hut, then the Marwadi for a thatched shelter, then say \emph{thā̃bho} again
\end{itemize}

\begin{tcolorbox}[breakable,colback=teal!4,colframe=teal!35!black,title={Before you move on}]
What does \mw{थांभो} mean? (\textquotedblleft{}A pillar\textquotedblright{}.) Say it once more.
\end{tcolorbox}
\section[\texorpdfstring{\mw{ईंट} (ī̃ṭ)}{ईंट (ī̃ṭ)}]{\mw{ईंट} (ī̃ṭ) — a brick}

\label{lesson:MW-C189-it}



\begin{quote}
Before the new one: say the Marwadi for a thatched shelter, then the Marwadi for a pillar.
\end{quote}

\subsection*{You'll want to know: \mw{ईंट}}
\textbf{\mw{ईंट}} — \emph{ī̃ṭ} — \textquotedblleft{}a brick\textquotedblright{}.

\begin{quote}
\mw{ईंट} \mw{लाल} \mw{है।}
\end{quote}

\begin{quote}
\emph{ī̃ṭ lāl hai.}
\end{quote}

\begin{quote}
The brick is red.
\end{quote}

A feminine noun.

\begin{cousinweb}
\textbf{\mw{ईंट}} (\emph{ī̃ṭ}) is inherited from Sanskrit \emph{iṣṭakā}, \textquotedblleft{}a brick\textquotedblright{}; CDIAL lists Hindi \emph{ī̃ṭ} (CDIAL 1600). CDIAL gives no Marwari form in this entry.
\end{cousinweb}

\begin{grammarlens}[title={What you've built}]
One more word for this chapter.
\end{grammarlens}

\subsection*{Your turn}
\begin{itemize}
\raggedright
  \item \emph{Say it:} \emph{ī̃ṭ}
  \item \emph{Say it:} \emph{ī̃ṭ}, once more
  \item \emph{Say it:} \emph{ī̃ṭ lāl hai.} — The brick is red.
  \item \emph{Recall:} say the Marwadi for a thatched shelter, then the Marwadi for a pillar, then say \emph{ī̃ṭ} again
\end{itemize}

\begin{tcolorbox}[breakable,colback=teal!4,colframe=teal!35!black,title={Before you move on}]
What does \mw{ईंट} mean? (\textquotedblleft{}A brick\textquotedblright{}.) Say it once more.
\end{tcolorbox}
\section[\texorpdfstring{\mw{चूनो} (chūno)}{चूनो (chūno)}]{\mw{चूनो} (chūno) — lime (for whitewash)}

\label{lesson:MW-C189-chuno}



\begin{quote}
Before the new one: say the Marwadi for a pillar, then the Marwadi for a brick.
\end{quote}

\subsection*{You'll want to know: \mw{चूनो}}
\textbf{\mw{चूनो}} — \emph{chūno} — \textquotedblleft{}lime (for whitewash)\textquotedblright{}.

\begin{quote}
\mw{भींत} \mw{माथै} \mw{नुवो} \mw{चूनो} \mw{है।}
\end{quote}

\begin{quote}
\emph{bhī̃t māthai nuvo chūno hai.}
\end{quote}

\begin{quote}
There's fresh whitewash on the wall.
\end{quote}

A masculine noun.

\begin{cousinweb}
\textbf{\mw{चूनो}} (\emph{chūno}) is inherited from Sanskrit \emph{cūrṇa}, \textquotedblleft{}powder\textquotedblright{}; CDIAL lists Hindi \emph{cūnā}, \textquotedblleft{}lime\textquotedblright{}, and an Old Marwari form meaning \textquotedblleft{}flour\textquotedblright{} (CDIAL 4889).
\end{cousinweb}

\begin{grammarlens}[title={What you've built}]
One more word for this chapter.
\end{grammarlens}

\subsection*{Your turn}
\begin{itemize}
\raggedright
  \item \emph{Say it:} \emph{chūno}
  \item \emph{Say it:} \emph{chūno}, once more
  \item \emph{Say it:} \emph{bhī̃t māthai nuvo chūno hai.} — There's fresh whitewash on the wall.
  \item \emph{Recall:} say the Marwadi for a pillar, then the Marwadi for a brick, then say \emph{chūno} again
\end{itemize}

\begin{tcolorbox}[breakable,colback=teal!4,colframe=teal!35!black,title={Before you move on}]
What does \mw{चूनो} mean? (\textquotedblleft{}Lime (for whitewash)\textquotedblright{}.) Say it once more.
\end{tcolorbox}
\section[\texorpdfstring{\mw{देहळी} (dehḷī)}{देहळी (dehḷī)}]{\mw{देहळी} (dehḷī) — a threshold}

\label{lesson:MW-C189-dehli}



\begin{quote}
Before the new one: say the Marwadi for a brick, then the Marwadi for whitewash.
\end{quote}

\subsection*{You'll want to know: \mw{देहळी}}
\textbf{\mw{देहळी}} — \emph{dehḷī} — \textquotedblleft{}a threshold\textquotedblright{}.

\begin{quote}
\mw{छोरो} \mw{देहळी} \mw{माथै} \mw{बैठै} \mw{है।}
\end{quote}

\begin{quote}
\emph{chhoro dehḷī māthai baiṭhai hai.}
\end{quote}

\begin{quote}
The boy sits on the threshold.
\end{quote}

A feminine noun.

\begin{cousinweb}
\textbf{\mw{देहळी}} (\emph{dehḷī}) is inherited from Sanskrit \emph{dehalī}, \textquotedblleft{}a threshold\textquotedblright{}; CDIAL lists Hindi \emph{dehlī} (CDIAL 6559). CDIAL gives no Marwari form in this entry.
\end{cousinweb}

\begin{grammarlens}[title={What you've built}]
One more word for this chapter.
\end{grammarlens}

\subsection*{Your turn}
\begin{itemize}
\raggedright
  \item \emph{Say it:} \emph{dehḷī}
  \item \emph{Say it:} \emph{dehḷī}, once more
  \item \emph{Say it:} \emph{chhoro dehḷī māthai baiṭhai hai.} — The boy sits on the threshold.
  \item \emph{Recall:} say the Marwadi for a brick, then the Marwadi for whitewash, then say \emph{dehḷī} again
\end{itemize}

\begin{tcolorbox}[breakable,colback=teal!4,colframe=teal!35!black,title={Before you move on}]
What does \mw{देहळी} mean? (\textquotedblleft{}A threshold\textquotedblright{}.) Say it once more.
\end{tcolorbox}
\section[\texorpdfstring{\mw{चौबारो} (chaubāro)}{चौबारो (chaubāro)}]{\mw{चौबारो} (chaubāro) — an open-sided room; a room open on all sides}

\label{lesson:MW-C189-chaubaro}



\begin{quote}
Before the new one: say the Marwadi for whitewash, then the Marwadi for a threshold.
\end{quote}

\subsection*{You'll want to know: \mw{चौबारो}}
\textbf{\mw{चौबारो}} — \emph{chaubāro} — \textquotedblleft{}an open-sided room; a room open on all sides\textquotedblright{}.

\begin{quote}
\mw{चौबारो} \mw{बड़ो} \mw{है।}
\end{quote}

\begin{quote}
\emph{chaubāro baṛo hai.}
\end{quote}

\begin{quote}
The open room is big.
\end{quote}

A masculine noun.

\begin{cousinweb}
\textbf{\mw{चौबारो}} (\emph{chaubāro}) is built from Sanskrit \emph{catur}, \textquotedblleft{}four\textquotedblright{}, and \emph{dvāra}, \textquotedblleft{}door\textquotedblright{}; CDIAL lists Hindi \emph{caubārā} and Sindhi \emph{caubāro}, a room or house open on four sides (CDIAL 4611). CDIAL gives no Marwari form in this entry.
\end{cousinweb}

\begin{grammarlens}[title={What you've built}]
One more word for this chapter.
\end{grammarlens}

\subsection*{Your turn}
\begin{itemize}
\raggedright
  \item \emph{Say it:} \emph{chaubāro}
  \item \emph{Say it:} \emph{chaubāro}, once more
  \item \emph{Say it:} \emph{chaubāro baṛo hai.} — The open room is big.
  \item \emph{Recall:} say the Marwadi for whitewash, then the Marwadi for a threshold, then say \emph{chaubāro} again
\end{itemize}

\begin{tcolorbox}[breakable,colback=teal!4,colframe=teal!35!black,title={Before you move on}]
What does \mw{चौबारो} mean? (\textquotedblleft{}An open-sided room; a room open on all sides\textquotedblright{}.) Say it once more.
\end{tcolorbox}
